% GENERATED FILE. Edit canonical lessons, then run npm run generate:books.
% canonical-source-hash: fnv1a64:6dd4f789881259e4
% canonical-lessons: PT-C526-episodio, PT-C526-temporada, PT-C526-telenovela, PT-C526-legendado, PT-C526-dobragem

\chapter{Séries e telenovelas}
\label{ch:pt-series-e-telenovelas}
\hlchaptermodality{526}

\begin{chapteropening}
\textbf{By the end of this chapter:} \emph{I can talk about series, seasons, subtitles and dubbing.}

Complete the last lesson of chapter 526: I can talk about series, seasons, subtitles and dubbing.
\end{chapteropening}

\section[o episódio]{o episódio — episode}

\label{lesson:PT-C526-episodio}



\begin{quote}
Before the new one: say the Portuguese for translation, then the Portuguese for translator.
\end{quote}

\subsection*{You'll want to know: o episódio}
\textbf{o episódio} — \textquotedblleft{}episode\textquotedblright{}.

\begin{quote}
Vi três episódios seguidos.
\end{quote}

\begin{quote}
\emph{I watched three episodes in a row.}
\end{quote}

Also an incident: \emph{um episódio estranho}.

\begin{cousinweb}
\emph{episódio} comes from French \emph{épisode} (EtymDB), from Greek \emph{epeisódion} \textquotedblleft{}an added part\textquotedblright{}, built on \emph{hodós} \textquotedblleft{}a way\textquotedblright{}.
\end{cousinweb}

\begin{grammarlens}[title={What you've built}]
One more word for telling, explaining and describing at B1.
\end{grammarlens}

\subsection*{Your turn}
\begin{itemize}
\raggedright
  \item \emph{Say it:} \emph{o episódio}
  \item \emph{Say it:} \emph{o episódio}, once more
  \item \emph{Say it:} a sentence of your own with \emph{o episódio}
  \item \emph{Recall:} say the Portuguese for translation, then the Portuguese for translator, then say \emph{o episódio} again
\end{itemize}

\begin{tcolorbox}[breakable,colback=teal!4,colframe=teal!35!black,title={Before you move on}]
What does o episódio mean? (\textquotedblleft{}Episode\textquotedblright{}.) Say it once more.
\end{tcolorbox}
\section[a temporada]{a temporada — season (of a series)}

\label{lesson:PT-C526-temporada}



\begin{quote}
Before the new one: say the Portuguese for translator, then the Portuguese for episode.
\end{quote}

\subsection*{You'll want to know: a temporada}
\textbf{a temporada} — \textquotedblleft{}season (of a series)\textquotedblright{}.

\begin{quote}
A segunda temporada é melhor do que a primeira.
\end{quote}

\begin{quote}
\emph{The second season is better than the first.}
\end{quote}

Also a season of the year for work or holidays: \emph{a temporada alta}.

\begin{cousinweb}
\emph{temporada} is built on \emph{tempo} \textquotedblleft{}time\textquotedblright{}, which is inherited from Latin \emph{tempus} (Lewis \& Short).
\end{cousinweb}

\begin{grammarlens}[title={What you've built}]
One more word for telling, explaining and describing at B1.
\end{grammarlens}

\subsection*{Your turn}
\begin{itemize}
\raggedright
  \item \emph{Say it:} \emph{a temporada}
  \item \emph{Say it:} \emph{a temporada}, once more
  \item \emph{Say it:} a sentence of your own with \emph{a temporada}
  \item \emph{Recall:} say the Portuguese for translator, then the Portuguese for episode, then say \emph{a temporada} again
\end{itemize}

\begin{tcolorbox}[breakable,colback=teal!4,colframe=teal!35!black,title={Before you move on}]
What does a temporada mean? (\textquotedblleft{}Season (of a series)\textquotedblright{}.) Say it once more.
\end{tcolorbox}
\section[a telenovela]{a telenovela — soap opera}

\label{lesson:PT-C526-telenovela}



\begin{quote}
Before the new one: say the Portuguese for episode, then the Portuguese for season.
\end{quote}

\subsection*{You'll want to know: a telenovela}
\textbf{a telenovela} — \textquotedblleft{}soap opera\textquotedblright{}.

\begin{quote}
A minha avó não perde a telenovela das nove.
\end{quote}

\begin{quote}
\emph{My grandmother never misses the nine o'clock soap.}
\end{quote}

Often shortened to \emph{a novela}.

\begin{cousinweb}
\emph{telenovela} is built from \emph{tele-}, from Greek \emph{tēle} \textquotedblleft{}far\textquotedblright{}, and \emph{novela}, from Italian \emph{novella} \textquotedblleft{}a tale\textquotedblright{} (EtymDB).
\end{cousinweb}

\begin{grammarlens}[title={What you've built}]
One more word for telling, explaining and describing at B1.
\end{grammarlens}

\subsection*{Your turn}
\begin{itemize}
\raggedright
  \item \emph{Say it:} \emph{a telenovela}
  \item \emph{Say it:} \emph{a telenovela}, once more
  \item \emph{Say it:} a sentence of your own with \emph{a telenovela}
  \item \emph{Recall:} say the Portuguese for episode, then the Portuguese for season, then say \emph{a telenovela} again
\end{itemize}

\begin{tcolorbox}[breakable,colback=teal!4,colframe=teal!35!black,title={Before you move on}]
What does a telenovela mean? (\textquotedblleft{}Soap opera\textquotedblright{}.) Say it once more.
\end{tcolorbox}
\section[legendado]{legendado — subtitled}

\label{lesson:PT-C526-legendado}



\begin{quote}
Before the new one: say the Portuguese for season, then the Portuguese for soap opera.
\end{quote}

\subsection*{You'll want to know: legendado}
\textbf{legendado} — \textquotedblleft{}subtitled\textquotedblright{}.

\begin{quote}
Prefiro ver os filmes legendados.
\end{quote}

\begin{quote}
\emph{I prefer to watch films with subtitles.}
\end{quote}

Subtitles are \emph{as legendas}.

\begin{cousinweb}
\emph{legendado} is built on \emph{legenda} \textquotedblleft{}a caption, a subtitle\textquotedblright{}, from Latin \emph{legenda} \textquotedblleft{}things to be read\textquotedblright{} (EtymDB), from \emph{legere} \textquotedblleft{}to read\textquotedblright{}.
\end{cousinweb}

\begin{grammarlens}[title={What you've built}]
One more word for telling, explaining and describing at B1.
\end{grammarlens}

\subsection*{Your turn}
\begin{itemize}
\raggedright
  \item \emph{Say it:} \emph{legendado}
  \item \emph{Say it:} \emph{legendado}, once more
  \item \emph{Say it:} a sentence of your own with \emph{legendado}
  \item \emph{Recall:} say the Portuguese for season, then the Portuguese for soap opera, then say \emph{legendado} again
\end{itemize}

\begin{tcolorbox}[breakable,colback=teal!4,colframe=teal!35!black,title={Before you move on}]
What does legendado mean? (\textquotedblleft{}Subtitled\textquotedblright{}.) Say it once more.
\end{tcolorbox}
\section[a dobragem]{a dobragem — dubbing}

\label{lesson:PT-C526-dobragem}



\begin{quote}
Before the new one: say the Portuguese for soap opera, then the Portuguese for subtitled.
\end{quote}

\subsection*{You'll want to know: a dobragem}
\textbf{a dobragem} — \textquotedblleft{}dubbing\textquotedblright{}.

\begin{quote}
Os desenhos animados para crianças têm dobragem.
\end{quote}

\begin{quote}
\emph{Children's cartoons are dubbed.}
\end{quote}

In Brazil: \emph{a dublagem}. In Portugal most films for adults are subtitled.

\begin{cousinweb}
\emph{dobragem} is built on \emph{dobrar} \textquotedblleft{}to fold, to double, to dub\textquotedblright{}, which goes back to Late Latin \emph{duplāre} \textquotedblleft{}to double\textquotedblright{} (EtymDB), from \emph{duplus}.
\end{cousinweb}

\begin{grammarlens}[title={What you've built}]
One more word for telling, explaining and describing at B1.
\end{grammarlens}

\subsection*{Your turn}
\begin{itemize}
\raggedright
  \item \emph{Say it:} \emph{a dobragem}
  \item \emph{Say it:} \emph{a dobragem}, once more
  \item \emph{Say it:} a sentence of your own with \emph{a dobragem}
  \item \emph{Recall:} say the Portuguese for soap opera, then the Portuguese for subtitled, then say \emph{a dobragem} again
\end{itemize}

\begin{tcolorbox}[breakable,colback=teal!4,colframe=teal!35!black,title={Before you move on}]
What does a dobragem mean? (\textquotedblleft{}Dubbing\textquotedblright{}.) Say it once more.
\end{tcolorbox}
